\documentclass[12pt]{amsart}
%   \overfullrule=5pt
\usepackage[active]{srcltx}
\usepackage{calc,amssymb,amsthm,amsmath,amscd, eucal,ulem}
\usepackage{alltt}
%\usepackage[left=1in,top=1in,right=1in,bottom=1in]{geometry}
\synctex=1
%\usepackage{mathtools}
\RequirePackage[dvipsnames,usenames]{color}
\let\mffont=\sf
\normalem
\input{mabliautoref.sty}
\input{xy}
\xyoption{all}
\input{kmacros3.sty}
\usepackage{tikz}
\usepackage{amsfonts, mathrsfs}
\usepackage{calligra}
\usepackage{stmaryrd} %power series bracket
%\usepackage{dcpic, pictexwd}
%\usepackage{pdfsync}
\usepackage[left=1.1in,top=1in,right=1.1in,bottom=1in]{geometry}
\usepackage{enumitem}

\numberwithin{equation}{theorem}

%\renewcommand{\thefootnote}{\fnsymbol{footnote}}
\newcommand{\D}{\displaystyle}
\newcommand{\til}{\widetilde}
\newcommand{\ol}{\overline}
\newcommand{\F}{\mathbb{F}}
\DeclareMathOperator{\E}{E}
\renewcommand{\:}{\colon}
\newcommand{\eg}{{\itshape e.g.} }
\renewcommand{\m}{\mathfrak{m}}
\renewcommand{\n}{\mathfrak{n}}
\newcommand{\Tt}{{\mathfrak{T}}}
\newcommand{\calC}{\mathcal{C}}
\newcommand{\RamiT}{\mathcal{R_{T}}}
\DeclareMathOperator{\edim}{edim}
\DeclareMathOperator{\length}{length}
\DeclareMathOperator{\monomials}{monomials}
\DeclareMathOperator{\Fitt}{Fitt}
\DeclareMathOperator{\depth}{depth}
\newcommand{\Frob}[2]{{#1}^{1/p^{#2}}}
\newcommand{\Frobp}[2]{{(#1)}^{1/p^{#2}}}
\newcommand{\FrobP}[2]{{\left(#1\right)}^{1/p^{#2}}}
\newcommand{\extends}{extends over $\Tt${}}
\newcommand{\extension}{extension over $\Tt${}}
\newcommand{\fg}{\pi_1^{\textnormal{\'{e}t}}} %Fundamental group
\DeclareMathOperator{\ns}{ns}
%\newcommand{\Spec}{\text{Spec}} %Spectrum
\newcommand{\pSpec}{\text{Spec}^{\circ}} %Puntured Spectrum
\newcommand{\etInCdOne}{\etale{} in codimension $1$}
\DeclareMathOperator{\DIV}{Div}
%\renewcommand{\char}{\charact}


\usepackage{setspace}
\usepackage{hyperref}
%\usepackage{enumerate}
\usepackage{graphicx}
\usepackage[all,cmtip]{xy}

\usepackage{verbatim}

\theoremstyle{theorem}
\newtheorem*{mainthm}{Main Theorem}
\newtheorem*{mainthma}{Main Theorem A}
\newtheorem*{mainthmb}{Main Theorem B}
\newtheorem*{mainthmc}{Main Theorem C}
\newtheorem*{atheorem}{Theorem}


\newcommand{\sol}[1]{ \vskip 6pt{\color{blue} {\bf Solution:} #1} \vskip 6pt}
%The todo box!
\def\todo#1{\textcolor{Mahogany}%
{\footnotesize\newline{\color{Mahogany}\fbox{\parbox{\textwidth-15pt}{\textbf{todo:
} #1}}}\newline}}

%What is going on?  Why isn't \O a script O?!?
\renewcommand{\O}{\mathscr O}

\begin{document}
\title{Exercises for Characteristic $p$ commutative algebra\\April 3rd, 2017}
\author{Due, Friday April 14th, 2017}
%\address{Department of Mathematics\\ University of Utah\\ Salt Lake City\\ UT 84112}
%\email{schwede@math.utah.edu}


%  \subjclass[2010]{14F18, 13A35, 14B05}

\maketitle

%\section{Reduction to characteristic $p > 0$}
%\chapter{Frobenius and Kunz's theorem}
%\bibliographystyle{skalpha}
%\bibliography{MainBib}

\begin{enumerate}[label=(\arabic*)]
\item Suppose $R$ is an $F$-finite normal domain and that $0 \neq \phi \in \Hom_R(F^e_* R, R)$ and $0 \neq \psi \in \Hom_R(F^d_* R, R)$.  Find a formula for $\Delta_{\phi \circ F^e_* \psi}$ in terms of $\Delta_{\phi}$ and $\Delta_{\psi}$.
\item Suppose that $\Delta \geq 0$ is a $\bQ$-divisor on a normal domain $R$.  Consider the subset
\[
S := \Hom_R(F^e_* R(\lceil (p^e - 1)\Delta \rceil), R) \subseteq \Hom_R(F^e_* R, R)
\]
induced by contravariance of $\Hom$.

Show that $S$ is equal to:
\[
\begin{array}{rl}
= & \{ \phi \in \Hom_R(F^e_* R, R)\;|\; \phi \text{ factors through the inclusion $R \subseteq R(\lceil (p^e-1)\Delta \rceil)$} \} \\
= & \{ \phi \in \Hom_R(F^e_* R, R) \;|\; \Delta_{\phi} \geq \Delta \}
\end{array}
\]
\item Suppose that $L / K$ is a finite extension of characteristic $p > 0$ fields and $x \in L \setminus K$ but $x^p \in K$.  Show that if $\phi : K^{1/p^e} \to K$ extends to $L^{1/p^e} \to L$, then $\phi$ is the zero map on $K$.
\item Using the previous problem, conclude in general that if $L/K$ is finite and inseparable, then no nonzero $\phi : K^{1/p^e} \to K$ can extend to $L^{1/p^e} \to L$.
    \vskip 3pt
    \emph{Hint:} Suppose such a map did extend to $\phi_L$.  Let $K'$ denote the separable closure of $K$ in $L$. Show that $\phi_L(K'^{1/p^e}) \subseteq K'$ and the reduce to the previous exercise.
\item Now fix a normal $F$-finite Noetherian domain $R$ (with $\Hom_R(F_* R, \omega_R) \cong F_* \omega_R$ if you'd like) and a nonzero $\phi \in \Hom_R(F^e_* R, R)$.  In this case we define the \emph{test ideal of $(R, \phi)$} to be the smallest nonzero ideal $J$ such that
    \[
    \phi(F^e_* J) \subseteq J,
    \]
    if it exists.  The test ideal is denoted by $\tau(R, \phi)$.
    In the following steps, we will show that it always exists.
    \begin{enumerate}
    \item Suppose first that $R$ is strongly $F$-regular and that $\phi$ generates $\Hom_R(F^e_* R, R)$ as an $F^e_* R$-module.  Show that for every $0 \neq c \in R$, there is an $n > 0$ such that $\phi^n(F^{ne}_* cR) = R$.  Conclude that $\tau(R, \phi) = R$.
        \vskip 3pt
        \emph{Hint:} We know there exists a $d > 0$ and $\psi \in \Hom_R(F^d_* R, R)$ such that $\psi(F^e_* c) = 1$.  Show that you can choose $d$ such that $e$ divides $d$. In this case, $\psi$ is a pre-multiple of some $\phi^n$.
    \item Show there exists $b \in R$ such that $R_b = R[b^{-1}]$ is strongly $F$-regular and $\phi_b$, the image of $\phi$, generates $\Hom_{R_b}(F^e_* R_b, R_b)$ as an $F^e_* R_b$-module.
        \vskip 3pt
        \emph{Hint:} First invert some $b'$ to make $R_{b'}$ regular (and hence strongly $F$-regular).  Then invert $b''$ contained in all the ideals that make up $D_{\phi}$.  Let $b = b'b''$.  Show that the image of $\phi$ in $\Hom_{R_{b}}(F^e_* R_b, R_b)$ corresponds to the zero divisor.
    \item Now fix $b$ as in the previous problem.  For any $0 \neq c \in R$, show that $b^{m_c} \in \phi^n(F^{ne}_* cR)$ for some $n > 0$ ($m_c$ depends on $c$).
        \vskip 3pt
        \emph{Hint:} We know that $1 \in \phi^n_b(F^{ne}_* cR_b)$ for some $n$.  Clear denominators.
    \item We can write $b^{m_1} \in \phi^{n_1}(F^{n_1 e}_* R)$ for some integers $n_1, m_1$.  Show that $b^{2 m_1} = (b^{m_1})^2 \in \phi^{l n_1}(F^{ln_1e}_* R)$ for all integers $l$.
        \vskip 3pt
        \emph{Hint:} Start by showing that $b^{\lceil m_1/p^{n_1 e} \rceil} b^{m_1} \in \phi^{n_1}(F^{n_1e}_* b^{m_1}R) \subseteq \phi^{2n_1}(F^{2n_1e}_* R)$.  Then do the general case.
        %\phi^{n}(F^{ne}_* b^{\lceil m_1/p^{ne} \rceil + m_1} R)
    \item Now fix $0 \neq c \in R$, and so $b^{m_c} \in \phi^{n_c}(F^{n_c e}_* c R)$ for some integers $m_c$ and $n_c$.  Conclude that $b^{3 m_1} \in \phi^N(F^{Ne}_* cR)$ for some integer $N > 0$.
        \vskip 3pt
        \emph{Hint:}  Apply $\phi^{n_1}$ until the result is obtained.
    \item Set $a = b^{3m_1}$.  Suppose that $J$ is any nonzero ideal such that $\phi(F^e_* J) \subseteq J$.  Show that $a \in J$.
    \item Show that $\tau(R, \phi) = \sum_{n \geq 0} \phi^n(F^{ne}_* aR)$ and hence the test ideal exists.
    \end{enumerate}
    \vskip 24pt
    \noindent
    {\bf Remark.}  Remember, $\phi$ corresponds to a divisor $\Delta_{\phi}$.  One can define $\tau(R, \Delta_{\phi}) := \tau(R, \phi)$, note that any two $\phi$ defining the same divisor are unit multiples.  Now, start with $R$ of finite type over a field of characteristic zero with some divisor $\Delta$ such that $K_R + \Delta \sim_{\bQ} 0$, then for all $p \gg 0$ we have that $(1-p^e)(K_R + \Delta) \sim 0$ for some $e$.  Hence if one reduces\footnote{In the case that $R$ is finite type over $\bQ$, we write that $R = R_{\bZ} \otimes_{\bZ} \bQ$ for some $R_{\bZ}$ of finite type over $\bQ$, then $R_p = R_{\bZ} \otimes_{\bZ} {\bZ/p\bZ}$.} $R$ to $R_p$, a characteristic $p \gg 0$ ring, one can choose $\phi_p$ corresponding to $\Delta_p$ (the reduction mod $p$ of $\Delta$ -- done for the primes defining $\Delta$).  In this case, it is a theorem that if $\mJ(R, \Delta)$ is the multiplier ideal, then
    \[
        \mJ(R, \Delta) \;\mathrm{mod}\; p = \tau(R_p, \Delta_p)
    \]
    for $p \gg 0$.  This was proven by Takagi in this generality.  Smith, Hara and Hara-Yoshida previously or simultaneously proved similar results.
\end{enumerate}



\end{document}


